\providecommand{\chEightCross}{%
\begin{tikzpicture}[font=\small,>=Stealth]
\draw[dashed,dcrred] (0,0) ellipse (3.0 and 2.8);
\node[text=dcrred] at (0,3.05) {可接触人群};
\node[align=center] at (0,1.95) {符合纳入标准的样本};
\draw[->] (1.7,1.9) to[bend left] (2,1.3);
\draw[dcrred] (-2.7,-1.3) rectangle (2.7,1.3);
\node[text=dcrpurple,align=center] at (0,.7) {测量当前变量\\获取既往资料（如相关）};
\node[draw,align=center] (x) at (-1.35,-.55) {符合\\排除标准};
\node[draw,align=center] (y) at (1.35,-.55) {不同意\\参加};
\node (e) at (0,-2) {排除};\draw[->] (e)--(x);\draw[->] (e)--(y);
\end{tikzpicture}}
\providecommand{\chEightCohort}[1]{%
\begin{tikzpicture}[font=\footnotesize,>=Stealth,x=.9cm,y=.9cm]
\draw[dashed,dcrred] (0,0) ellipse (2.4 and 2.2);
\node[text=dcrred] at (0,2.5) {可接触人群};
\node[align=center] at (0,1.6) {符合选择标准的样本};
\draw[->] (1.1,1.6) to[bend left] (1.5,1.05);
\draw[dcrred] (-2.1,-1.05) rectangle (2.1,1.05);
\node[text=dcrpurple,align=center] at (0,0) {测量基线变量\\[9pt]\ifnum#1=0 获取既往资料（如相关）\else 保存标本（可选）\fi};
\draw[thick,dcrteal,->] (2.1,.8)--(9.1,.8);
\node[text=dcrteal,align=center] at (5.6,1.15) {随访参与者};
\draw[thick,dcrteal,->] (2.1,-.8)--(9.1,-.8);
\ifnum#1>0
\ifnum#1=2
\foreach \x in {3.05,5.45,7.7}{\node[text=dcrgold,align=center,text width=9mm,inner sep=1pt,font=\fontsize{8}{9.5}\selectfont] at (\x,0) {结局\\发生时\\测量};}
\else
\node[text=dcrgold,align=center] at (5.7,0) {结局发生时测量};
\fi
\foreach \x in {3.2,5.8,8}{\draw[->,dcrteal] (\x,-.8) to[bend left=30] (\x+.6,-2.05);}
\node[text=dcrteal] at (6,-2.4) {失访};
\else
\draw[->,dcrteal] (6,-.8) to[bend left=35] (7,-1.8);
\node[text=dcrteal] at (6,-2.3) {极少失访或无人失访};
\fi
\ifnum#1=2
\foreach \x in {4.1,6.8}{\node[draw=dcrpurple,fill=white,text=dcrpurple,rotate=90] at (\x,0) {更新预测因素};}
\fi
\node[draw=dcrgold,text=dcrgold,align=center] at (9.6,0) {\ifnum#1=0 随访结束时\\测量结局\else 可选择在随访\\结束时测量结局\fi};
\draw[dcrteal,thick,->] (2.1,-3)--(9,-3);
\node[fill=white,text=dcrteal] at (5.5,-3) {随访时间\ifnum#1=0 （短）\fi};
\draw[dcrteal] (-2.1,-1.05)--(-2.1,-2.7)--(2.1,-2.7)--(2.1,-1.05);
\node[text=dcrteal] at (0,-3) {研究起始};
\end{tikzpicture}}
\providecommand{\chEightDouble}{%
\begin{tikzpicture}[font=\footnotesize,>=Stealth,x=.9cm,y=.9cm]
\foreach \yy in {0,-5.5}{
\begin{scope}[yshift=\yy cm]
\draw[dashed,dcrred] (0,0) ellipse (2.4 and 2.1);
\node[align=center] at (0,1.6) {符合选择标准的样本};
\draw[->] (1.1,1.6) to[bend left] (1.5,1);
\draw[dcrred] (-2.1,-1) rectangle (2.1,1);
\node[text=dcrpurple,align=center] at (0,0) {测量基线变量\\保存标本（很少）};
\draw[thick,dcrteal,->] (2.1,.7)--(9,.7);
\draw[thick,dcrteal,->] (2.1,-.7)--(9,-.7);
\node[text=dcrteal] at (5.4,1) {随访参与者};
\node[text=dcrgold] at (5.6,.1) {结局发生时测量};
\foreach \x in {3.3,6,8}{\draw[->,dcrteal] (\x,-.7) to[bend left=30] (\x+.6,-1.9);}
\node[text=dcrteal] at (6,-2.3) {失访};
\node[draw=dcrgold,text=dcrgold,align=center] at (9.7,0) {可选择在随访\\结束时测量结局};
\end{scope}}
\node[text=dcrred,align=center] at (0,2.65) {可接触的暴露人群\\（人数常有限）};
\node[text=dcrred,align=center] at (0,-2.8) {可接触的未暴露或不同暴露人群\\（通常较大）};
\node[fill=white,align=center] at (0,-4.25) {符合选择标准的样本\\（人数比可为1:1、2:1等）};
\draw[dcrteal,thick,->] (2.1,-8.8)--(9,-8.8);
\node[fill=white,text=dcrteal] at (5.4,-8.8) {随访时间};
\node[text=dcrteal] at (0,-8.8) {研究起始};
\end{tikzpicture}}
